\chapter{Configura\c{c}\~oes para executar a app}

\section{Vis\~ao geral}

Para executar o projeto BaseBallX s\~ao necess\'arios tr\^es elementos principais:

\begin{itemize}[leftmargin=*]
  \item um ambiente Python com as depend\^encias da aplica\c{c}\~ao;
  \item um reposit\'orio \texttt{baseball} ativo no GraphDB;
  \item a aplica\c{c}\~ao Django a correr localmente.
\end{itemize}

O fluxo de execu\c{c}\~ao pode ser resumido da seguinte forma:

\begin{enumerate}[leftmargin=*]
  \item criar e ativar o ambiente virtual Python;
  \item instalar depend\^encias;
  \item gerar o ficheiro RDF \texttt{rdf/baseball.n3};
  \item iniciar o GraphDB e carregar o reposit\'orio \texttt{baseball};
  \item aplicar as migrations do Django;
  \item iniciar o servidor web.
\end{enumerate}

\section{Pr\'e-requisitos}

Antes da execu\c{c}\~ao, o sistema deve ter instalado:

\begin{itemize}[leftmargin=*]
  \item \textbf{Python 3};
  \item \textbf{GraphDB Desktop}
  \item \textbf{curl}, usado pelos scripts auxiliares;
  \item acesso ao diret\'orio do projeto com os dados na pasta \texttt{archive/}.
\end{itemize}

Garantir que a porta \texttt{7200} est\'a livre para o GraphDB 
\begin{center}
\texttt{http://localhost:7200}
\end{center}

\section{Execu\c{c}\~ao manual}

\subsection{1. Criar o ambiente virtual}

Na raiz do projeto:

\begin{verbatim}
python3 -m venv venv
source venv/bin/activate
pip install -r requirements.txt
\end{verbatim}

\subsection{2. Gerar o RDF}

Se o ficheiro RDF ainda n\~ao existir deve ser executada a convers\~ao:

\begin{verbatim}
python3 rdf/convert.py
\end{verbatim}

Este comando gera o ficheiro:

\begin{center}
\texttt{rdf/baseball.n3}
\end{center}

\subsection{3. Iniciar o GraphDB e criar o reposit\'orio}

Existem duas formas de preparar o GraphDB.

\textbf{Op\c{c}\~ao A: usar o script auxiliar}

\begin{verbatim}
./scripts/start_graphdb.sh
\end{verbatim}

Este script:

\begin{itemize}[leftmargin=*]
  \item verifica se o GraphDB est\'a acess\'ivel;
  \item tenta abrir o GraphDB Desktop se o servi\c{c}o ainda n\~ao estiver ativo;
  \item cria o reposit\'orio \texttt{baseball} se este ainda n\~ao existir;
  \item importa automaticamente o ficheiro \texttt{rdf/baseball.n3}.
\end{itemize}

Por omiss\~ao, o script usa:

\begin{itemize}[leftmargin=*]
  \item \texttt{GRAPHDB\_URL=http://localhost:7200}
  \item \texttt{REPOSITORY\_ID=baseball}
  \item \texttt{RDF\_FILE=\$ROOT\_DIR/rdf/baseball.n3}
  \item \texttt{GRAPHDB\_DESKTOP\_BIN=/opt/graphdb-desktop/bin/graphdb-desktop}
\end{itemize}

% \textbf{Op\c{c}\~ao B: fazer manualmente no GraphDB Desktop}
% 
% \begin{enumerate}[leftmargin=*]
%   \item abrir o GraphDB Desktop;
%   \item iniciar o servi\c{c}o local;
%   \item criar um reposit\'orio com o nome \texttt{baseball};
%   \item importar o ficheiro \texttt{rdf/baseball.n3} para esse reposit\'orio.
% \end{enumerate}
% 
% \subsection{4. Preparar a base de dados Django}

A aplica\c{c}\~ao usa SQLite para persistir contas, sugest\~oes, sess\~oes e resultados do quiz. Depois de ativar o ambiente virtual, devem ser aplicadas as migrations:

\begin{verbatim}
cd webapp
python manage.py migrate
\end{verbatim}

\subsection{5. Iniciar o servidor web}

Ainda na pasta \texttt{webapp}, iniciar a aplica\c{c}\~ao Django:

\begin{verbatim}
python manage.py runserver
\end{verbatim}

Depois disso, a aplica\c{c}\~ao fica dispon\'ivel em:

\begin{center}
\texttt{http://127.0.0.1:8000/}
\end{center}

\section{Fluxo completo recomendado}

O fluxo mais direto para executar o projeto desde o in\'icio \'{e}:

\begin{verbatim}
python3 -m venv venv
source venv/bin/activate
pip install -r requirements.txt
python3 rdf/convert.py
./scripts/start_graphdb.sh
cd webapp
python manage.py migrate
python manage.py runserver
\end{verbatim}

\section{Parar e reconstruir o reposit\'orio}

Se for necess\'ario apagar o reposit\'orio do GraphDB e voltar a import\'a-lo, existe um script dedicado:

\begin{verbatim}
./scripts/delete_graphdb_repo.sh
\end{verbatim}

Depois da remo\c{c}\~ao, basta correr novamente:

\begin{verbatim}
./scripts/start_graphdb.sh
\end{verbatim}

Isto volta a criar o reposit\'orio \texttt{baseball} e a importar o ficheiro RDF.
